% Generated. Do not edit.
% Projection scriptural-date-location/section, version 1: the Scriptural Date and
% Location page, generated whole from research/chronology.toml, the Scripture
% chronology corpus and the source library by tools/tpt document-types refresh.
% document: liturgy/roman-rite/postconciliar/roman-missal-third-edition-en-us-2011/propers/temporal/pc-s52-twenty-sixth-sunday-in-ordinary-time-year-a
\runninghead{Scriptural Date and Location}
\section{Scriptural Date and Location}
\label{sec:date-location}

{\small Each passage the Mass reads or sings from Scripture, in canonical order. The Context column names what is dated --- the event the passage narrates, a setting it remembers, the traditional attribution of its author, or the composition of its book --- and where it belongs; the Date column gives the answer of the Scripture chronology corpus for that subject, with every alternative it records. Texts of the Mass that are not drawn from Scripture have no row. The evidence for each date follows the table.\par}
\medskip
% Generated. Do not edit.
% Deterministic projection of research/chronology.toml.
% schema: 5
% document: liturgy/roman-rite/postconciliar/roman-missal-third-edition-en-us-2011/propers/temporal/pc-s52-twenty-sixth-sunday-in-ordinary-time-year-a
% generated-by: tools/tpt proper-chronology annotations
\newcommand{\chronologyannotationclaim}[6]{#6}
\newcommand{\chronologyannotationreach}[3]{}
\newcommand{\chronologyannotationgroup}[3]{#3}
\newcommand{\chronologyannotation}[1]{%
  \ifcsname triptychchronologyannotation@#1\endcsname
    \csname triptychchronologyannotation@#1\endcsname
  \else
    \PackageError{triptych}{No chronology annotation for #1}{}%
  \fi}
% appointment-dependency: guidance/liturgy/postconciliar-propers-registry.md sha256=376860a7c37a4c75bb6d880cd1b260183be1877a07e3cd01660a9e2854e2ea72
% appointment-dependency: src/claude/liturgy/roman-rite/postconciliar/roman-missal-third-edition-en-us-2011/propers/registry/README.md sha256=22626b14e0506b69c0ddc254b0ade2b6f7c9ae7cf8ef98fe356746a8f47fb1eb
% appointment-dependency: src/claude/liturgy/roman-rite/postconciliar/roman-missal-third-edition-en-us-2011/propers/registry/formula-dispositions.md sha256=37735c7edc0fbf692a691c7c4b0487b80831108daf91621b91f7b8858e190879
% appointment-dependency: src/claude/liturgy/roman-rite/postconciliar/roman-missal-third-edition-en-us-2011/propers/temporal/pc-s52-twenty-sixth-sunday-in-ordinary-time-year-a/instance/manifest.md sha256=92e3872be41c139bc810df3c15a555aa94a34c88c5a0d685c88ba6c944896afd
% appointment-dependency: src/claude/liturgy/roman-rite/postconciliar/roman-missal-third-edition-en-us-2011/propers/temporal/pc-s52-twenty-sixth-sunday-in-ordinary-time-year-a/propers/verified.md sha256=17e8a79f4c0a8572a14f3948d6976cc7e19c68f284e447f400d827429e5ae3ba
% appointment-dependency: src/claude/liturgy/roman-rite/postconciliar/roman-missal-third-edition-en-us-2011/propers/temporal/pc-s52-twenty-sixth-sunday-in-ordinary-time-year-a/research/chronology-inputs.toml sha256=7ec09fd99e1fd72b263235bae19118e250623352e692d8f74cba3c4a728a8bad
% appointment-dependency: src/claude/liturgy/roman-rite/postconciliar/roman-missal-third-edition-en-us-2011/propers/temporal/shared/ordinary-time/weeks/26/propers/verified.md sha256=fa7cff717694e821885255d19ee7f01fddd03d59d2e2a4374afb96ddd55ee376
% appointment-note: entrance-antiphon: Daniel 3:31, 29, 30, 43, 42; adaptation; vulgate numbering; canonical verse queries, not appointed wording.
% appointment-note: collect: Composed oration with no scriptural locator; its vocabulary of sparing and mercy is a loose echo of no one verse, which the page-2 rule excludes.
% appointment-note: first-reading: Ezekiel 18:25-28; appointed; vulgate numbering; canonical verse queries, not appointed wording.
% appointment-note: responsorial-psalm: Psalm 25:4-5, 6-7, 8-9; appointed; hebrew numbering; canonical verse queries, not appointed wording.
% appointment-note: responsorial-psalm: Psalm 25:6a; response; hebrew numbering; whole-verse query envelope, not subverse chronology or appointed wording.
% appointment-note: second-reading-long: Philippians 2:1-11; appointed; vulgate numbering; canonical verse queries, not appointed wording.
% appointment-note: second-reading-short: Philippians 2:1-5; appointed; vulgate numbering; canonical verse queries, not appointed wording.
% appointment-note: gospel-acclamation: John 10:27; appointed; vulgate numbering; canonical verse queries, not appointed wording.
% appointment-note: gospel: Matthew 21:28-32; appointed; vulgate numbering; canonical verse queries, not appointed wording.
% appointment-note: prayer-over-offerings: Composed oration with no scriptural locator.
% appointment-note: communion-antiphon-a: Cf. Psalm 118:49-50; adaptation; vulgate numbering; canonical verse queries, not appointed wording.
% appointment-note: communion-antiphon-b: 1 John 3:16; appointed; vulgate numbering; canonical verse queries, not appointed wording.
% appointment-note: prayer-after-communion: Composed oration with no scriptural locator. The words coheredes and compatimur also stand together in Romans 8:17, and mortem annuntiando recalls 1 Corinthians 11:26; these are loose echoes the Missal does not cite, which the page-2 rule excludes.
% entrance-antiphon: Entrance Antiphon
\expandafter\def\csname triptychchronologyannotation@entrance-antiphon\endcsname{%
  \chronologyannotationgroup{historical-setting}{preferred}{\textbf{Historical setting}: \chronologyannotationreach{Dan.3.29}{Dan}{true}\chronologyannotationreach{Dan.3.30}{Dan}{true}\chronologyannotationreach{Dan.3.31}{Dan}{true}\chronologyannotationreach{Dan.3.42}{Dan}{true}\chronologyannotationreach{Dan.3.43}{Dan}{true}\chronologyannotationclaim{israel.exile.seventy-years}{historical-setting}{catholic-traditional-v1}{preferred}{seventy years}{The seventy years of servitude to Babylon, ``seventy years'' (duration)}.}%
  \space \chronologyannotationgroup{composition}{preferred}{\textbf{Composition}: \chronologyannotationreach{Dan.3.29}{Dan.3.24-90}{false}\chronologyannotationreach{Dan.3.30}{Dan.3.24-90}{false}\chronologyannotationreach{Dan.3.31}{Dan.3.24-90}{false}\chronologyannotationreach{Dan.3.42}{Dan.3.24-90}{false}\chronologyannotationreach{Dan.3.43}{Dan.3.24-90}{false}\chronologyannotationclaim{composition.book-of-daniel.prayer-of-azarias}{composition}{catholic-traditional-v1}{preferred}{most likely long after the Exile}{Most likely long after the Exile}.}%
}
% first-reading: First Reading
\expandafter\def\csname triptychchronologyannotation@first-reading\endcsname{%
  \chronologyannotationgroup{historical-setting}{preferred}{\textbf{Historical setting}: Preferred: \chronologyannotationreach{Ezech.18.25}{Ezech}{true}\chronologyannotationreach{Ezech.18.26}{Ezech}{true}\chronologyannotationreach{Ezech.18.27}{Ezech}{true}\chronologyannotationreach{Ezech.18.28}{Ezech}{true}\chronologyannotationclaim{israel.exile.second-captivity}{historical-setting}{catholic-traditional-v1}{preferred}{in the eighth year of his reign}{The second captivity of Juda, under Jechonias, in the eighth year of his reign}; alternatives: \chronologyannotationreach{Ezech.18.25}{Ezech}{true}\chronologyannotationreach{Ezech.18.26}{Ezech}{true}\chronologyannotationreach{Ezech.18.27}{Ezech}{true}\chronologyannotationreach{Ezech.18.28}{Ezech}{true}\chronologyannotationclaim{israel.exile.second-captivity}{historical-setting}{catholic-traditional-v1}{alternate}{A.M. 3405}{The second captivity of Juda, under Jechonias, A.M. 3405}, \chronologyannotationreach{Ezech.18.25}{Ezech}{true}\chronologyannotationreach{Ezech.18.26}{Ezech}{true}\chronologyannotationreach{Ezech.18.27}{Ezech}{true}\chronologyannotationreach{Ezech.18.28}{Ezech}{true}\chronologyannotationclaim{israel.exile.second-captivity}{historical-setting}{catholic-traditional-v1}{alternate}{597 B.C.}{The second captivity of Juda, under Jechonias, B.C. 597}, \chronologyannotationreach{Ezech.18.25}{Ezech}{true}\chronologyannotationreach{Ezech.18.26}{Ezech}{true}\chronologyannotationreach{Ezech.18.27}{Ezech}{true}\chronologyannotationreach{Ezech.18.28}{Ezech}{true}\chronologyannotationclaim{israel.exile.second-captivity}{historical-setting}{catholic-traditional-v1}{alternate}{in the year 598 B.C.}{The second captivity of Juda, under Jechonias, B.C. 598}.}%
  \space \chronologyannotationgroup{prophecy-given}{preferred}{\textbf{Prophecy given}: \chronologyannotationreach{Ezech.18.25}{Ezech}{true}\chronologyannotationreach{Ezech.18.26}{Ezech}{true}\chronologyannotationreach{Ezech.18.27}{Ezech}{true}\chronologyannotationreach{Ezech.18.28}{Ezech}{true}\chronologyannotationclaim{israel.exile.ezechiel.ministry}{prophecy-given}{catholic-traditional-v1}{preferred}{he prophesied during at least twenty-two years}{The prophetic ministry of Ezechiel among the exiles by the river Chobar, ``he prophesied during at least twenty-two years'' (duration)}.}%
}
% responsorial-psalm: Responsorial Psalm
\expandafter\def\csname triptychchronologyannotation@responsorial-psalm\endcsname{%
  \chronologyannotationgroup{historical-setting}{research-pending}{\textbf{Historical setting} -- Occasion date unresolved.}%
  \space \chronologyannotationgroup{composition}{preferred}{\textbf{Composition}: \chronologyannotationreach{Ps.24.4}{Ps}{true}\chronologyannotationreach{Ps.24.5}{Ps}{true}\chronologyannotationreach{Ps.24.6}{Ps}{true}\chronologyannotationreach{Ps.24.7}{Ps}{true}\chronologyannotationreach{Ps.24.8}{Ps}{true}\chronologyannotationreach{Ps.24.9}{Ps}{true}\chronologyannotationclaim{critical.psalms.latest-composition-boundary}{composition}{catholic-critical-v1}{preferred}{before the Maccabean period, around 165 B.C.; no individual psalm can be dated securely}{Before c. 165 B.C.}}%
}
% second-reading-long: Second Reading (longer form)
\expandafter\def\csname triptychchronologyannotation@second-reading-long\endcsname{%
  \chronologyannotationgroup{composition}{disputed}{\textbf{Composition} -- disputed: \chronologyannotationreach{Phil.2.1}{Phil}{true}\chronologyannotationreach{Phil.2.10}{Phil}{true}\chronologyannotationreach{Phil.2.11}{Phil}{true}\chronologyannotationreach{Phil.2.2}{Phil}{true}\chronologyannotationreach{Phil.2.3}{Phil}{true}\chronologyannotationreach{Phil.2.4}{Phil}{true}\chronologyannotationreach{Phil.2.5}{Phil}{true}\chronologyannotationreach{Phil.2.6}{Phil}{true}\chronologyannotationreach{Phil.2.7}{Phil}{true}\chronologyannotationreach{Phil.2.8}{Phil}{true}\chronologyannotationreach{Phil.2.9}{Phil}{true}\chronologyannotationclaim{composition.epistle-to-the-philippians}{composition}{catholic-traditional-v1}{disputed}{(Philemon; Colossians; Ephesians; Philippians), 61}{A.D. 61}; \chronologyannotationreach{Phil.2.1}{Phil}{true}\chronologyannotationreach{Phil.2.10}{Phil}{true}\chronologyannotationreach{Phil.2.11}{Phil}{true}\chronologyannotationreach{Phil.2.2}{Phil}{true}\chronologyannotationreach{Phil.2.3}{Phil}{true}\chronologyannotationreach{Phil.2.4}{Phil}{true}\chronologyannotationreach{Phil.2.5}{Phil}{true}\chronologyannotationreach{Phil.2.6}{Phil}{true}\chronologyannotationreach{Phil.2.7}{Phil}{true}\chronologyannotationreach{Phil.2.8}{Phil}{true}\chronologyannotationreach{Phil.2.9}{Phil}{true}\chronologyannotationclaim{composition.epistle-to-the-philippians}{composition}{catholic-traditional-v1}{disputed}{at Rome (A.D. 62-64)}{A.D. 62--64}.}%
}
% second-reading-short: Second Reading (shorter form)
\expandafter\def\csname triptychchronologyannotation@second-reading-short\endcsname{%
  \chronologyannotationgroup{composition}{disputed}{\textbf{Composition} -- disputed: \chronologyannotationreach{Phil.2.1}{Phil}{true}\chronologyannotationreach{Phil.2.2}{Phil}{true}\chronologyannotationreach{Phil.2.3}{Phil}{true}\chronologyannotationreach{Phil.2.4}{Phil}{true}\chronologyannotationreach{Phil.2.5}{Phil}{true}\chronologyannotationclaim{composition.epistle-to-the-philippians}{composition}{catholic-traditional-v1}{disputed}{(Philemon; Colossians; Ephesians; Philippians), 61}{A.D. 61}; \chronologyannotationreach{Phil.2.1}{Phil}{true}\chronologyannotationreach{Phil.2.2}{Phil}{true}\chronologyannotationreach{Phil.2.3}{Phil}{true}\chronologyannotationreach{Phil.2.4}{Phil}{true}\chronologyannotationreach{Phil.2.5}{Phil}{true}\chronologyannotationclaim{composition.epistle-to-the-philippians}{composition}{catholic-traditional-v1}{disputed}{at Rome (A.D. 62-64)}{A.D. 62--64}.}%
}
% gospel-acclamation: Alleluia Verse
\expandafter\def\csname triptychchronologyannotation@gospel-acclamation\endcsname{%
  \chronologyannotationgroup{composition}{disputed}{\textbf{Composition} -- disputed: \chronologyannotationreach{John.10.27}{John}{true}\chronologyannotationclaim{composition.gospel-of-john}{composition}{catholic-traditional-v1}{disputed}{from the year 90 to 100 (approximately)}{c. A.D. 90--100}; \chronologyannotationreach{John.10.27}{John}{true}\chronologyannotationclaim{composition.gospel-of-john}{composition}{catholic-traditional-v1}{disputed}{96 or one of the succeeding years}{A.D. 96--100}.}%
}
% gospel: Gospel
\expandafter\def\csname triptychchronologyannotation@gospel\endcsname{%
  \chronologyannotationgroup{narrated-event}{research-pending}{\textbf{Event} -- No narrated-event date in the chronology corpus.}%
  \space \chronologyannotationgroup{composition}{disputed}{\textbf{Composition} -- disputed: \chronologyannotationreach{Matt.21.28}{Matt}{true}\chronologyannotationreach{Matt.21.29}{Matt}{true}\chronologyannotationreach{Matt.21.30}{Matt}{true}\chronologyannotationreach{Matt.21.31}{Matt}{true}\chronologyannotationreach{Matt.21.32}{Matt}{true}\chronologyannotationclaim{composition.gospel-of-matthew}{composition}{catholic-traditional-v1}{disputed}{about A.D. 38-45}{c. A.D. 38--45}; \chronologyannotationreach{Matt.21.28}{Matt}{true}\chronologyannotationreach{Matt.21.29}{Matt}{true}\chronologyannotationreach{Matt.21.30}{Matt}{true}\chronologyannotationreach{Matt.21.31}{Matt}{true}\chronologyannotationreach{Matt.21.32}{Matt}{true}\chronologyannotationclaim{composition.gospel-of-matthew}{composition}{catholic-traditional-v1}{disputed}{about the year 40-42}{c. A.D. 40--42}; \chronologyannotationreach{Matt.21.28}{Matt}{true}\chronologyannotationreach{Matt.21.29}{Matt}{true}\chronologyannotationreach{Matt.21.30}{Matt}{true}\chronologyannotationreach{Matt.21.31}{Matt}{true}\chronologyannotationreach{Matt.21.32}{Matt}{true}\chronologyannotationclaim{composition.gospel-of-matthew}{composition}{catholic-traditional-v1}{disputed}{the years 40-45}{A.D. 40--45}; \chronologyannotationreach{Matt.21.28}{Matt}{true}\chronologyannotationreach{Matt.21.29}{Matt}{true}\chronologyannotationreach{Matt.21.30}{Matt}{true}\chronologyannotationreach{Matt.21.31}{Matt}{true}\chronologyannotationreach{Matt.21.32}{Matt}{true}\chronologyannotationclaim{composition.gospel-of-matthew}{composition}{catholic-traditional-v1}{disputed}{about the year 60-68}{c. A.D. 60--68}; \chronologyannotationreach{Matt.21.28}{Matt}{true}\chronologyannotationreach{Matt.21.29}{Matt}{true}\chronologyannotationreach{Matt.21.30}{Matt}{true}\chronologyannotationreach{Matt.21.31}{Matt}{true}\chronologyannotationreach{Matt.21.32}{Matt}{true}\chronologyannotationclaim{composition.gospel-of-matthew}{composition}{catholic-traditional-v1}{disputed}{about the years 64-67}{c. A.D. 64--67}; \chronologyannotationreach{Matt.21.28}{Matt}{true}\chronologyannotationreach{Matt.21.29}{Matt}{true}\chronologyannotationreach{Matt.21.30}{Matt}{true}\chronologyannotationreach{Matt.21.31}{Matt}{true}\chronologyannotationreach{Matt.21.32}{Matt}{true}\chronologyannotationclaim{composition.gospel-of-matthew}{composition}{catholic-traditional-v1}{disputed}{about the year 50}{c. A.D. 50}.}%
}
% communion-antiphon-a: Communion Antiphon (first)
\expandafter\def\csname triptychchronologyannotation@communion-antiphon-a\endcsname{%
  \chronologyannotationgroup{traditional-attribution}{disputed}{\textbf{Traditional attribution} -- disputed: \chronologyannotationreach{Ps.118.49}{Ps.118}{true}\chronologyannotationreach{Ps.118.50}{Ps.118}{true}\chronologyannotationclaim{israel.monarchy.david-traditional-era}{traditional-attribution}{catholic-traditional-v1}{disputed}{reigned from 1055 to 1015 B.C.}{David (reign in the usual chronology), B.C. 1055--1015}.}%
  \space \chronologyannotationgroup{historical-setting}{research-pending}{\textbf{Historical setting} -- Occasion date unresolved.}%
  \space \chronologyannotationgroup{composition}{preferred}{\textbf{Composition}: \chronologyannotationreach{Ps.118.49}{Ps}{true}\chronologyannotationreach{Ps.118.50}{Ps}{true}\chronologyannotationclaim{critical.psalms.latest-composition-boundary}{composition}{catholic-critical-v1}{preferred}{before the Maccabean period, around 165 B.C.; no individual psalm can be dated securely}{Before c. 165 B.C.}}%
}
% communion-antiphon-b: Communion Antiphon (alternative)
\expandafter\def\csname triptychchronologyannotation@communion-antiphon-b\endcsname{%
  \chronologyannotationgroup{composition}{preferred}{\textbf{Composition}: \chronologyannotationreach{1John.3.16}{1John}{true}\chronologyannotationclaim{composition.first-epistle-of-st-john}{composition}{catholic-traditional-v1}{preferred}{from the year 90 to 100 (approximately)}{c. A.D. 90--100}.}%
}
% presentation: orientation-v3
\newcommand{\chronologyorientationrow}[2]{%
  \ifcsname triptychchronologyorientation@#1\endcsname
    \csname triptychchronologyorientation@#1\endcsname{#2}%
  \else
    \PackageError{triptych}{No chronology orientation row for #1}{}%
  \fi}
\expandafter\def\csname triptychchronologyorientation@row-2684d5253d65bb6c\endcsname#1{%
Responsorial Psalm: Ps. 25:4--5, 6--7, 8--9 (Heb.; Vg. 24), Ps. 25:6a (Heb.; Vg. 24) & Historical setting & Occasion date unresolved\\*[.1em]%
}
\expandafter\def\csname triptychchronologyorientation@row-bf73bd4c86f3b7be\endcsname#1{%
 & Composition: Psalms & \chronologyannotationreach{Ps.24.4}{Ps}{true}\chronologyannotationreach{Ps.24.5}{Ps}{true}\chronologyannotationreach{Ps.24.6}{Ps}{true}\chronologyannotationreach{Ps.24.7}{Ps}{true}\chronologyannotationreach{Ps.24.8}{Ps}{true}\chronologyannotationreach{Ps.24.9}{Ps}{true}\chronologyannotationclaim{critical.psalms.latest-composition-boundary}{composition}{catholic-critical-v1}{preferred}{before the Maccabean period, around 165 B.C.; no individual psalm can be dated securely}{Before c.~165 B.C.}\\[.3em]%
}
\expandafter\def\csname triptychchronologyorientation@row-551cbc6a6787985d\endcsname#1{%
Communion Antiphon (first): Cf. Psalm 118:49--50 (Vg.; Heb. 119) & Traditional attribution: David's reign; Hebron, then Jerusalem & Disputed: \chronologyannotationreach{Ps.118.49}{Ps.118}{true}\chronologyannotationreach{Ps.118.50}{Ps.118}{true}\chronologyannotationclaim{israel.monarchy.david-traditional-era}{traditional-attribution}{catholic-traditional-v1}{disputed}{reigned from 1055 to 1015 B.C.}{B.C.~1055--1015}\\*[.1em]%
}
\expandafter\def\csname triptychchronologyorientation@row-8f762950b4f3652e\endcsname#1{%
 & Historical setting & Occasion date unresolved\\*[.1em]%
}
\expandafter\def\csname triptychchronologyorientation@row-253466ade46ac58c\endcsname#1{%
 & Composition: Psalms & \chronologyannotationreach{Ps.118.49}{Ps}{true}\chronologyannotationreach{Ps.118.50}{Ps}{true}\chronologyannotationclaim{critical.psalms.latest-composition-boundary}{composition}{catholic-critical-v1}{preferred}{before the Maccabean period, around 165 B.C.; no individual psalm can be dated securely}{Before c.~165 B.C.}\\[.3em]%
}
\expandafter\def\csname triptychchronologyorientation@row-c488125d2bc454fe\endcsname#1{%
First Reading: Ezech. 18:25--28 & Historical setting & Date pending its source qualification\\*[.1em]%
}
\expandafter\def\csname triptychchronologyorientation@row-8567e6166a24b18b\endcsname#1{%
 & Prophecy given: The prophetic ministry of Ezechiel among the exiles by the river Chobar & \chronologyannotationreach{Ezech.18.25}{Ezech}{true}\chronologyannotationreach{Ezech.18.26}{Ezech}{true}\chronologyannotationreach{Ezech.18.27}{Ezech}{true}\chronologyannotationreach{Ezech.18.28}{Ezech}{true}\chronologyannotationclaim{israel.exile.ezechiel.ministry}{prophecy-given}{catholic-traditional-v1}{preferred}{he prophesied during at least twenty-two years}{``he prophesied during at least twenty-two years'' (duration)}\\[.3em]%
}
\expandafter\def\csname triptychchronologyorientation@row-93d43c952988f909\endcsname#1{%
Entrance Antiphon: Dan. 3:31, 29, 30, 43, 42 & Historical setting: The seventy years of servitude to Babylon; Judah and Babylon & \chronologyannotationreach{Dan.3.29}{Dan}{true}\chronologyannotationreach{Dan.3.30}{Dan}{true}\chronologyannotationreach{Dan.3.31}{Dan}{true}\chronologyannotationreach{Dan.3.42}{Dan}{true}\chronologyannotationreach{Dan.3.43}{Dan}{true}\chronologyannotationclaim{israel.exile.seventy-years}{historical-setting}{catholic-traditional-v1}{preferred}{seventy years}{``seventy years'' (duration)}\\*[.1em]%
}
\expandafter\def\csname triptychchronologyorientation@row-8375fb7b245fcec7\endcsname#1{%
 & Composition: Prayer of Azarias and Song of the Three Children & \chronologyannotationreach{Dan.3.29}{Dan.3.24-90}{false}\chronologyannotationreach{Dan.3.30}{Dan.3.24-90}{false}\chronologyannotationreach{Dan.3.31}{Dan.3.24-90}{false}\chronologyannotationreach{Dan.3.42}{Dan.3.24-90}{false}\chronologyannotationreach{Dan.3.43}{Dan.3.24-90}{false}\chronologyannotationclaim{composition.book-of-daniel.prayer-of-azarias}{composition}{catholic-traditional-v1}{preferred}{most likely long after the Exile}{Most likely long after the Exile}\\[.3em]%
}
\expandafter\def\csname triptychchronologyorientation@row-d2121e4192fb2d88\endcsname#1{%
Gospel: Matt. 21:28--32 & Narrated event & Date unresolved\\*[.1em]%
}
\expandafter\def\csname triptychchronologyorientation@row-1e127cc3406e8e75\endcsname#1{%
 & Composition: Matthew; Semitic original associated with Palestine & Disputed: \chronologyannotationreach{Matt.21.28}{Matt}{true}\chronologyannotationreach{Matt.21.29}{Matt}{true}\chronologyannotationreach{Matt.21.30}{Matt}{true}\chronologyannotationreach{Matt.21.31}{Matt}{true}\chronologyannotationreach{Matt.21.32}{Matt}{true}\chronologyannotationclaim{composition.gospel-of-matthew}{composition}{catholic-traditional-v1}{disputed}{about A.D. 38-45}{c.~A.D.~38--45}; \chronologyannotationreach{Matt.21.28}{Matt}{true}\chronologyannotationreach{Matt.21.29}{Matt}{true}\chronologyannotationreach{Matt.21.30}{Matt}{true}\chronologyannotationreach{Matt.21.31}{Matt}{true}\chronologyannotationreach{Matt.21.32}{Matt}{true}\chronologyannotationclaim{composition.gospel-of-matthew}{composition}{catholic-traditional-v1}{disputed}{about the year 40-42}{c.~A.D.~40--42}; \chronologyannotationreach{Matt.21.28}{Matt}{true}\chronologyannotationreach{Matt.21.29}{Matt}{true}\chronologyannotationreach{Matt.21.30}{Matt}{true}\chronologyannotationreach{Matt.21.31}{Matt}{true}\chronologyannotationreach{Matt.21.32}{Matt}{true}\chronologyannotationclaim{composition.gospel-of-matthew}{composition}{catholic-traditional-v1}{disputed}{the years 40-45}{A.D.~40--45}; \chronologyannotationreach{Matt.21.28}{Matt}{true}\chronologyannotationreach{Matt.21.29}{Matt}{true}\chronologyannotationreach{Matt.21.30}{Matt}{true}\chronologyannotationreach{Matt.21.31}{Matt}{true}\chronologyannotationreach{Matt.21.32}{Matt}{true}\chronologyannotationclaim{composition.gospel-of-matthew}{composition}{catholic-traditional-v1}{disputed}{about the year 60-68}{c.~A.D.~60--68}; \chronologyannotationreach{Matt.21.28}{Matt}{true}\chronologyannotationreach{Matt.21.29}{Matt}{true}\chronologyannotationreach{Matt.21.30}{Matt}{true}\chronologyannotationreach{Matt.21.31}{Matt}{true}\chronologyannotationreach{Matt.21.32}{Matt}{true}\chronologyannotationclaim{composition.gospel-of-matthew}{composition}{catholic-traditional-v1}{disputed}{about the years 64-67}{c.~A.D.~64--67}; \chronologyannotationreach{Matt.21.28}{Matt}{true}\chronologyannotationreach{Matt.21.29}{Matt}{true}\chronologyannotationreach{Matt.21.30}{Matt}{true}\chronologyannotationreach{Matt.21.31}{Matt}{true}\chronologyannotationreach{Matt.21.32}{Matt}{true}\chronologyannotationclaim{composition.gospel-of-matthew}{composition}{catholic-traditional-v1}{disputed}{about the year 50}{c.~A.D.~50}\\[.3em]%
}
\expandafter\def\csname triptychchronologyorientation@row-21cf30d7ca478297\endcsname#1{%
Alleluia Verse: John 10:27 & Composition: John; Ephesus traditionally & Disputed: \chronologyannotationreach{John.10.27}{John}{true}\chronologyannotationclaim{composition.gospel-of-john}{composition}{catholic-traditional-v1}{disputed}{from the year 90 to 100 (approximately)}{c.~A.D.~90--100}; \chronologyannotationreach{John.10.27}{John}{true}\chronologyannotationclaim{composition.gospel-of-john}{composition}{catholic-traditional-v1}{disputed}{96 or one of the succeeding years}{A.D.~96--100}\\[.3em]%
}
\expandafter\def\csname triptychchronologyorientation@row-ce4292186d26d172\endcsname#1{%
Second Reading (longer form): Phil. 2:1--11 & Composition: Philippians; Rome to Philippi in the cited account & Disputed: \chronologyannotationreach{Phil.2.1}{Phil}{true}\chronologyannotationreach{Phil.2.10}{Phil}{true}\chronologyannotationreach{Phil.2.11}{Phil}{true}\chronologyannotationreach{Phil.2.2}{Phil}{true}\chronologyannotationreach{Phil.2.3}{Phil}{true}\chronologyannotationreach{Phil.2.4}{Phil}{true}\chronologyannotationreach{Phil.2.5}{Phil}{true}\chronologyannotationreach{Phil.2.6}{Phil}{true}\chronologyannotationreach{Phil.2.7}{Phil}{true}\chronologyannotationreach{Phil.2.8}{Phil}{true}\chronologyannotationreach{Phil.2.9}{Phil}{true}\chronologyannotationclaim{composition.epistle-to-the-philippians}{composition}{catholic-traditional-v1}{disputed}{(Philemon; Colossians; Ephesians; Philippians), 61}{A.D.~61}; \chronologyannotationreach{Phil.2.1}{Phil}{true}\chronologyannotationreach{Phil.2.10}{Phil}{true}\chronologyannotationreach{Phil.2.11}{Phil}{true}\chronologyannotationreach{Phil.2.2}{Phil}{true}\chronologyannotationreach{Phil.2.3}{Phil}{true}\chronologyannotationreach{Phil.2.4}{Phil}{true}\chronologyannotationreach{Phil.2.5}{Phil}{true}\chronologyannotationreach{Phil.2.6}{Phil}{true}\chronologyannotationreach{Phil.2.7}{Phil}{true}\chronologyannotationreach{Phil.2.8}{Phil}{true}\chronologyannotationreach{Phil.2.9}{Phil}{true}\chronologyannotationclaim{composition.epistle-to-the-philippians}{composition}{catholic-traditional-v1}{disputed}{at Rome (A.D. 62-64)}{A.D.~62--64}\\[.3em]%
}
\expandafter\def\csname triptychchronologyorientation@row-597b8a6a642339e3\endcsname#1{%
Second Reading (shorter form): Phil. 2:1--5 & Composition: Philippians; Rome to Philippi in the cited account & Disputed: \chronologyannotationreach{Phil.2.1}{Phil}{true}\chronologyannotationreach{Phil.2.2}{Phil}{true}\chronologyannotationreach{Phil.2.3}{Phil}{true}\chronologyannotationreach{Phil.2.4}{Phil}{true}\chronologyannotationreach{Phil.2.5}{Phil}{true}\chronologyannotationclaim{composition.epistle-to-the-philippians}{composition}{catholic-traditional-v1}{disputed}{(Philemon; Colossians; Ephesians; Philippians), 61}{A.D.~61}; \chronologyannotationreach{Phil.2.1}{Phil}{true}\chronologyannotationreach{Phil.2.2}{Phil}{true}\chronologyannotationreach{Phil.2.3}{Phil}{true}\chronologyannotationreach{Phil.2.4}{Phil}{true}\chronologyannotationreach{Phil.2.5}{Phil}{true}\chronologyannotationclaim{composition.epistle-to-the-philippians}{composition}{catholic-traditional-v1}{disputed}{at Rome (A.D. 62-64)}{A.D.~62--64}\\[.3em]%
}
\expandafter\def\csname triptychchronologyorientation@row-9c4564480a6fa781\endcsname#1{%
Communion Antiphon (alternative): 1 John. 3:16 & Composition: First John; Ephesus probable & \chronologyannotationreach{1John.3.16}{1John}{true}\chronologyannotationclaim{composition.first-epistle-of-st-john}{composition}{catholic-traditional-v1}{preferred}{from the year 90 to 100 (approximately)}{c.~A.D.~90--100}\\[.3em]%
}

% Generated. Do not edit.
% Complete orientation-v3 table; wording, dates and order belong to the corpus.
\begin{chronologysharedtable}
\chronologyorientationrow{row-2684d5253d65bb6c}{}
\chronologyorientationrow{row-bf73bd4c86f3b7be}{}
\midrule
\chronologyorientationrow{row-551cbc6a6787985d}{}
\chronologyorientationrow{row-8f762950b4f3652e}{}
\chronologyorientationrow{row-253466ade46ac58c}{}
\midrule
\chronologyorientationrow{row-c488125d2bc454fe}{}
\chronologyorientationrow{row-8567e6166a24b18b}{}
\midrule
\chronologyorientationrow{row-93d43c952988f909}{}
\chronologyorientationrow{row-8375fb7b245fcec7}{}
\midrule
\chronologyorientationrow{row-d2121e4192fb2d88}{}
\chronologyorientationrow{row-1e127cc3406e8e75}{}
\midrule
\chronologyorientationrow{row-21cf30d7ca478297}{}
\midrule
\chronologyorientationrow{row-ce4292186d26d172}{}
\midrule
\chronologyorientationrow{row-597b8a6a642339e3}{}
\midrule
\chronologyorientationrow{row-9c4564480a6fa781}{}
\end{chronologysharedtable}
{\footnotesize Psalm numbering: Vg. = Vulgate; Heb. = Hebrew.\par}


\subsection{Evidence for each date}
\label{sec:date-location-evidence}
{\small For each passage, as the Scriptural Date and Location table names it: what is dated, how the corpus weighs each answer, the source's own words, and the sources cited. A subject the corpus does not date is reported as such and given no date.\par}
{\small
\par\smallskip\noindent\textit{Responsorial Psalm: Ps. 25:4--5, 6--7, 8--9 (Heb.; Vg. 24), Ps. 25:6a (Heb.; Vg. 24)}\par\nobreak
\begin{itemize}[leftmargin=1.2em,itemsep=1pt,topsep=1pt,parsep=0pt]
\item \textbf{Historical setting}. No date: the chronology corpus has not yet dated this (research pending).
\item \textbf{Composition} (preferred). Psalms: Before c.~165 B.C. (``before the Maccabean period, around 165 B.C.; no individual psalm can be dated securely''; Catholic critical chronology; New American Bible, Revised Edition --- USCCB web state (2026), Psalms introduction).
\end{itemize}
\par\smallskip\noindent\textit{Communion Antiphon (first): Cf. Ps. 118:49--50 (Vg.; Heb. 119)}\par\nobreak
\begin{itemize}[leftmargin=1.2em,itemsep=1pt,topsep=1pt,parsep=0pt]
\item \textbf{Traditional attribution} (disputed). David's reign: B.C.~1055--1015 (``reigned from 1055 to 1015 B.C.''; The Catholic Encyclopedia, Volume 4 (New York, 1908), article 'David, King,' printed p. 642, opening paragraph). Place: Hebron, then Jerusalem (The Holy Bible, Douay-Rheims (Challoner revision), Project Gutenberg eBook 1581: 2 Kings 5:5).
\item \textbf{Historical setting}. As under Responsorial Psalm: Ps. 25:4--5, 6--7, 8--9 (Heb.; Vg. 24), Ps. 25:6a (Heb.; Vg. 24).
\item \textbf{Composition}. As under Responsorial Psalm: Ps. 25:4--5, 6--7, 8--9 (Heb.; Vg. 24), Ps. 25:6a (Heb.; Vg. 24).
\end{itemize}
\par\smallskip\noindent\textit{First Reading: Ezech. 18:25--28}\par\nobreak
\begin{itemize}[leftmargin=1.2em,itemsep=1pt,topsep=1pt,parsep=0pt]
\item \textbf{Historical setting}. No date: the corpus holds a date for this that it reports from another edition by exception, which is not printed until the corpus records the source's own qualification of it (research pending).
\item \textbf{Prophecy given} (preferred). The prophetic ministry of Ezechiel among the exiles by the river Chobar: ``he prophesied during at least twenty-two years'' (duration) (The Catholic Encyclopedia, Volume 5 (New York, 1909); The Holy Bible, Douay-Rheims (Challoner revision), Project Gutenberg eBook 1581: Ezech. 1:2, 29:17).
\end{itemize}
\par\smallskip\noindent\textit{Entrance Antiphon: Dan. 3:31, 29, 30, 43, 42}\par\nobreak
\begin{itemize}[leftmargin=1.2em,itemsep=1pt,topsep=1pt,parsep=0pt]
\item \textbf{Historical setting} (preferred). The seventy years of servitude to Babylon: ``seventy years'' (duration) (The Holy Bible, Douay-Rheims (Challoner revision), Project Gutenberg eBook 1581: Jer. 25:11, 25:12). Place: Judah and Babylon (The Holy Bible, Douay-Rheims (Challoner revision), Project Gutenberg eBook 1581: Jer. 25:11, 25:12).
\item \textbf{Composition} (preferred). Prayer of Azarias and Song of the Three Children: Most likely long after the Exile (The Catholic Encyclopedia, Volume 4 (New York, 1908)).
\end{itemize}
\par\smallskip\noindent\textit{Gospel: Matt. 21:28--32}\par\nobreak
\begin{itemize}[leftmargin=1.2em,itemsep=1pt,topsep=1pt,parsep=0pt]
\item \textbf{Narrated event}. No date: the chronology corpus has not yet dated this (research pending).
\item \textbf{Composition} (disputed). The Gospel of St Matthew: c.~A.D.~38--45 (``about A.D. 38-45''; The Catholic Encyclopedia, Volume 10 (1911); The Catholic Encyclopedia, Volume 14 (New York, 1912)); c.~A.D.~40--42 (``about the year 40-42''; The Catholic Encyclopedia, Volume 10 (1911)); A.D.~40--45 (``the years 40-45''; The Catholic Encyclopedia, Volume 10 (1911)); c.~A.D.~60--68 (``about the year 60-68''; The Catholic Encyclopedia, Volume 10 (1911)); c.~A.D.~64--67 (``about the years 64-67''; The Catholic Encyclopedia, Volume 10 (1911)); c.~A.D.~50 (``about the year 50''; The Catholic Encyclopedia, Volume 14 (New York, 1912)). Place: Semitic original associated with Palestine (The Catholic Encyclopedia, Volume 10 (1911)).
\end{itemize}
\par\smallskip\noindent\textit{Alleluia Verse: John 10:27}\par\nobreak
\begin{itemize}[leftmargin=1.2em,itemsep=1pt,topsep=1pt,parsep=0pt]
\item \textbf{Composition} (disputed). The Gospel of St John: c.~A.D.~90--100 (``from the year 90 to 100 (approximately)''; The Catholic Encyclopedia, Volume 14 (New York, 1912)); A.D.~96--100 (``96 or one of the succeeding years''; The Catholic Encyclopedia, Volume 8 (New York, 1910)). Place: Ephesus traditionally (The Catholic Encyclopedia, Volume 8 (New York, 1910)).
\end{itemize}
\par\smallskip\noindent\textit{Second Reading (longer form): Phil. 2:1--11}\par\nobreak
\begin{itemize}[leftmargin=1.2em,itemsep=1pt,topsep=1pt,parsep=0pt]
\item \textbf{Composition} (disputed). The Epistle to the Philippians: A.D.~61 (``(Philemon; Colossians; Ephesians; Philippians), 61''; The Catholic Encyclopedia, Volume 11 (New York, 1911)); A.D.~62--64 (``at Rome (A.D. 62-64)''; The Catholic Encyclopedia, Volume 12 (1911)). Place: Rome to Philippi in the cited account (The Catholic Encyclopedia, Volume 12 (1911); The Holy Bible, Douay-Rheims (Challoner revision), Project Gutenberg eBook 1581: Phil. 1:1).
\end{itemize}
\par\smallskip\noindent\textit{Second Reading (shorter form): Phil. 2:1--5}\par\nobreak
\begin{itemize}[leftmargin=1.2em,itemsep=1pt,topsep=1pt,parsep=0pt]
\item \textbf{Composition}. As under Second Reading (longer form): Phil. 2:1--11.
\end{itemize}
\par\smallskip\noindent\textit{Communion Antiphon (alternative): 1 John. 3:16}\par\nobreak
\begin{itemize}[leftmargin=1.2em,itemsep=1pt,topsep=1pt,parsep=0pt]
\item \textbf{Composition} (preferred). The First Epistle of St John: c.~A.D.~90--100 (``from the year 90 to 100 (approximately)''; The Catholic Encyclopedia, Volume 14 (New York, 1912); The Catholic Encyclopedia, Volume 8 (New York, 1910)). Place: Ephesus probable (The Catholic Encyclopedia, Volume 8 (New York, 1910)).
\end{itemize}
\par}
